\vspace{-0.3cm}
\section{晶圆级加速器设备模型}
\vspace{-0.2cm}

\subsection{PLMR 模型}
\vspace{-0.1cm}
我们提出 PLMR 模型，用于刻画晶圆级加速器独特的硬件属性，并由此推导利用这类新兴硬件所需满足的系统要求。

\begin{enumerate}[label=(\arabic*),leftmargin=0.5cm,noitemsep,topsep=0pt]
\item \textbf{海量\underline{并行性（P）}}：
晶圆级加速器可以容纳数百万个并行核心，而 GPU 通常只有数千个。每个核心都带有本地硬件流水线，可在周期粒度上重叠数据输入、数据输出、计算与内存访问。因此，计算必须进行超大规模划分，并通过细粒度调度来重叠计算、内存访问和 NoC 通信。

\item \textbf{高度非均匀的内存访问\underline{延迟（L）}}：在网格中跨核心访问内存会产生高度非均匀的延迟。对于一个按矩形排列、包含 $N_w \times N_h$ 个核心的网格，任意两核之间的最大 NoC 跳数为 $N_w+N_h$，最坏情况下的内存访问延迟为 $\alpha(N_w+N_h)+\beta r$，其中 $r<N_w+N_h$ 表示通信路径上的路由阶段数。这里，$\alpha$ 是单跳传输延迟，即消息在某个核心依据路由硬件预配置规则被直接转发时产生的成本，它随跳数增加；$\beta$ 是单次路由延迟，即消息在某个核心转发时由软件解析并重写报头所产生的开销~\cite{wafer-reduce}。通常 $\alpha<\beta$。在拥有一百万个核心的网格中，最大跳数与路由次数可达数千，从而使本地与远程内存访问的延迟相差最高一千倍。因此，尽量减少长距离通信对性能至关重要。

\item \textbf{受限的单核本地\underline{内存（M）}}：每个核心只有很小的本地内存（数十 KB 到数 MB），因为更大容量会降低性能与能效~\cite{sram-wiki}。因此，必须把计算数据显式划分为细粒度块，使其完整容纳于每个核心的本地内存限制之内。

\item \textbf{受限的\underline{路由资源（R）}}：集成数百万核心的晶圆级加速器会严格限制每个核心的路由电路复杂度或路由表大小。例如在 Cerebras WSE-2 上，每个核心只能识别带 5-bit 地址编码的消息报头。因此，每个核心最多支持 $2^5$ 条不同路由路径，软件系统必须仔细规划这些路径。对于长距离远程通信，一对核心可以消耗路由资源建立直接路由路径，只产生 $\alpha$ 延迟；但如果路由路径数超过硬件限制，消息就必须经由多个中间核心中继，从而引入额外的 $\beta$ 延迟。
\end{enumerate}

我们预计这些属性会持续成立，因为它们源自硬件及其制造工艺的基本特性。PLMR 模型既适用于当前的 Cerebras WSE，也适用于未来的 Tesla Dojo 等晶圆级设备。甚至一些采用网格 NoC、但并非晶圆级的设备，例如 Tenstorrent Blackhole~\cite{tenstorrent}，也可以用 PLMR 表示，只需调整并行度（P）、网格规模（L）等参数，或放宽本地内存（M）和路由资源（R）的约束。

\vspace{-3mm}
\subsection{最先进方法的局限}
\vspace{-1mm}

借助 PLMR 模型，我们分析现有 AI 系统为何无法充分利用晶圆级加速器。要在晶圆级加速器上运行 LLM，通常有两种选择：（i）把每个核心的分布式本地内存抽象成共享内存，并通过 NoC 直接访问放置在远程核心上的数据；（ii）把计算显式划分到分布式核心，并用消息传递交换必要数据。我们分析两类有代表性的系统：一类是面向 GPU 等共享内存架构的 LLM runtime 或 DNN 编译器，例如 Ladder~\cite{ladder}；另一类是面向分布式片上内存架构的 SOTA 编译器，例如用于 GraphCore IPU 的 T10~\cite{t10}。

\mypar{共享内存系统}
Ladder 等基于共享内存的 DNN 编译器通常假设底层内存层次中的访问模式是均匀的，因而无法容忍晶圆级加速器访问远程内存时高达数千倍的延迟差异（不满足 L）。此外，这些编译器~\cite{tvm,rammer,ansor,flextensor,roller,welder,ladder}常把重点放在计算划分上，而较少优化数据划分。这很容易造成大量数据重复，并违反内存约束要求（不满足 M）。最后，这类编译器不知道各核心之间的通信距离，因而也无法妥善处理路由资源约束。

\mypar{分布式内存系统} T10 系统~\cite{t10}面向带片上交叉开关的 AI 加速器设计；这种交叉开关保证访问同一芯片上任意其他核心内存的延迟恒定。T10 能处理小容量本地内存并平衡通信负载，因此满足内存约束（M）和路由资源限制（R）。但在 PLMR 设备上，它没有考虑变化的内存访问延迟（不满足 L），并且只能扩展到数千而非数百万核心（不满足 P）。
